\section{Empirical strategy}\label{sec:strategy}

For each documentary-reviewed revision $e$ at institution $i(e)$ with event year $E_e$, we form a stack containing the
revising institution and institutions with no threshold revision anywhere in the window $[E_e-4,E_e+5]$, observed over
the same calendar years. Pooling the stacks, we estimate
\begin{equation}
Y_{ist}=\sum_{k\neq-1}\delta_k D^k_{ist}
+\sum_{k\neq-1}\theta_k D^k_{ist}S_{e(s)}
+\gamma\log R_{it}+\alpha_{is}+\lambda_{stc(i)}+\varepsilon_{ist},
\label{eq:stacked}
\end{equation}
where $s$ indexes stacks, $D^k_{ist}$ indicates the revising institution $k$ years from its revision, and
$S_e\in\{-1,+1\}$ equals $+1$ for an upward revision and $-1$ for a downward revision. The
institution-by-stack effects $\alpha_{is}$ absorb level differences. The stack-by-year effects
$\lambda_{stc(i)}$ are interacted with research-size tercile and public/private status, and $R_{it}$ is research
expenditure. Year $-1$ is omitted. The stacked structure uses clean contemporaneous comparisons and avoids the
contaminated-comparison problems that arise in conventional two-way fixed-effects event studies under staggered timing
\citep{GoodmanBacon2021,CallawaySantAnna2021,SunAbraham2021,CengizEtAl2019}.

The paths $\delta_k+\theta_k$ and $\delta_k-\theta_k$ describe upward and downward revisions relative to their controls,
and $2\theta_k$ is the difference between the two. The headline estimand is
\begin{equation}
\Delta=\frac{1}{6}\sum_{k=0}^{5}2\theta_k.
\end{equation}
This difference-in-differences-in-direction comparison removes changes that accompany revision generally. It does not
remove changes that systematically accompany \emph{the direction} of revision. If universities that tighten their
policies are simultaneously reorganizing TTOs or changing budgets in ways that reduce licensing, $\Delta$ may capture
those changes as well. The identifying comparison therefore requires that, absent direction-specific changes, upward and
downward revisers would have evolved similarly relative to their clean controls. Pre-revision paths and balance
diagnostics are informative about this condition but cannot establish it.

Standard errors are clustered by institution. The effective variation in direction comes from 34 events, including only
eight upward revisions, so we place greater weight on randomization-style permutation inference than on asymptotic
clustered $p$-values. Holding event dates, stacks, outcomes, and controls fixed, we repeatedly reassign which eight of the
34 events are labeled upward and recompute $\Delta$. The preferred results use 100,000 random assignments preserving the
number of upward events; the Monte Carlo standard error of the licensing permutation $p$-value is approximately 0.0006.
This procedure tests exchangeability of direction labels under the null. Because direction was not randomly assigned,
exchangeability is an assumption, not a feature of the institutional design. We therefore report pre-revision balance,
pre-trend tests, leave-one-event-out estimates, and the stricter 29-event sample alongside the permutation results.
Benjamini--Hochberg $q$-values adjust the five main-outcome permutation tests as a family.

The annual panel remains supporting evidence. We regress each outcome on PCSI measured zero to five years earlier with
institution and year fixed effects and lagged log research expenditure, log licensing FTEs, and inventor royalty share,
clustering by institution. Because most annual PCSI values are carried forward and revision timing is less precise in the
annual representation, these regressions are not treated as the principal design.
